\documentclass[11pt,a4paper]{report}

% ============================================================================
% RetailVision Edge AI — Manual de Usuario
% Universidad Santo Tomás, Seccional Tunja
% ============================================================================

\usepackage[utf8]{inputenc}
\usepackage[T1]{fontenc}
\usepackage{lmodern}
\usepackage[spanish,es-noindentfirst]{babel}
\usepackage[a4paper,margin=2.5cm]{geometry}
\usepackage{xcolor}
\usepackage{graphicx}
\usepackage{tabularx}
\usepackage{longtable}
\usepackage{booktabs}
\usepackage{enumitem}
\usepackage{fancyhdr}
\usepackage{titlesec}
\usepackage{hyperref}
\usepackage{array}
\usepackage{colortbl}
\usepackage{microtype}
\sloppy

% --- Paleta (coherente con el dashboard) ------------------------------------
\definecolor{rvAccent}{HTML}{C76A00}
\definecolor{rvAccentDark}{HTML}{8A4600}
\definecolor{rvInk}{HTML}{2C3242}
\definecolor{rvMuted}{HTML}{5B6478}
\definecolor{rvLine}{HTML}{CCD4E2}
\definecolor{rvCodeBg}{HTML}{F2F3F6}
\definecolor{rvGreen}{HTML}{1F9D66}
\definecolor{rvRed}{HTML}{D33B3B}

\titleformat{\chapter}[display]
  {\normalfont\bfseries\color{rvInk}}
  {\color{rvAccent}\Large\chaptertitlename\ \thechapter}
  {6pt}{\Huge}
\titleformat{\section}
  {\normalfont\Large\bfseries\color{rvInk}}{\thesection}{0.6em}{}
\titleformat{\subsection}
  {\normalfont\large\bfseries\color{rvAccentDark}}{\thesubsection}{0.6em}{}

\setlength{\headheight}{14.5pt}
\addtolength{\topmargin}{-0.5pt}
\pagestyle{fancy}
\fancyhf{}
\renewcommand{\headrulewidth}{0.4pt}
\renewcommand{\footrulewidth}{0pt}
\fancyhead[L]{\footnotesize\color{rvMuted} RetailVision Edge AI}
\fancyhead[R]{\footnotesize\color{rvMuted} \leftmark}
\renewcommand{\chaptermark}[1]{\markboth{#1}{}}
\fancyfoot[C]{\small\color{rvMuted} \thepage}

\hypersetup{
  colorlinks=true,
  linkcolor=rvAccentDark,
  urlcolor=rvAccentDark,
  citecolor=rvAccentDark,
  pdftitle={RetailVision Edge AI -- Manual de Usuario},
  pdfauthor={Equipo RetailVision Edge AI -- Universidad Santo Tomas, Seccional Tunja},
}

\newcommand{\rvnota}[1]{%
  \par\noindent\fcolorbox{rvAccent}{rvCodeBg}{%
    \parbox{0.95\linewidth}{\small \textbf{\color{rvAccentDark}Nota:} #1}%
  }\par\vspace{4pt}
}
\newcommand{\rvconsejo}[1]{%
  \par\noindent\fcolorbox{rvGreen}{rvCodeBg}{%
    \parbox{0.95\linewidth}{\small \textbf{\color{rvGreen}Consejo:} #1}%
  }\par\vspace{4pt}
}
\newcommand{\rvpaso}[1]{\textbf{\color{rvAccentDark}#1}}

\setlist[itemize]{leftmargin=1.4em}
\setlist[enumerate]{leftmargin=1.6em}

\begin{document}

\begin{titlepage}
\centering
\vspace*{1.5cm}
{\Huge\bfseries\color{rvAccent} RetailVision Edge AI\par}
\vspace{0.4cm}
{\Large\color{rvInk} Manual de Usuario\par}
\vspace{0.3cm}
{\large\color{rvMuted} Guía de uso del panel de analítica de comportamiento en tienda\par}
\vspace{2.2cm}

{\large\color{rvInk}\textbf{Universidad Santo Tomás, Seccional Tunja}\par}
\vspace{0.2cm}
{\normalsize\color{rvMuted} Desarrollo académico grupal\par}
\vspace{0.15cm}
{\normalsize\color{rvMuted} Docente asesor: Sergio Arley Puerto Moreno\par}

\vspace{2.2cm}
{\normalsize\color{rvInk}\textbf{Equipo}\par}
\vspace{0.3cm}
\begin{tabular}{l}
Silvana Sofía Siza Soriano \\
Yeikov Jared Ramos Sánchez \\
Sarita Toloza Buitrago \\
Juan David Sandoval Bello \\
Thomas Felipe Cepeda Mogollón \\
Wendy Mariana Vargas Ramírez \\
Andrés Felipe Galeano Téllez \\
Sara Valentina Caicedo Duarte \\
\end{tabular}

\vfill
{\large\color{rvMuted} 2026\par}
\end{titlepage}

\tableofcontents
\clearpage

% ============================================================================
\chapter{Bienvenida}
% ============================================================================

\section{¿Qué es RetailVision Edge AI?}

RetailVision Edge AI es un panel web que analiza el video de las cámaras de una góndola o
estante y responde, de forma automática, tres preguntas que normalmente son difíciles de medir
en una tienda física:

\begin{itemize}
  \item ¿Cuántas personas se detienen frente a cada estante?
  \item ¿Quién \textbf{toma} un producto y lo \textbf{devuelve} sin comprarlo?
  \item ¿Qué zonas de la góndola funcionan mejor y cuáles pasan desapercibidas?
\end{itemize}

Con esa información, el sistema propone, en lenguaje sencillo, cómo redistribuir los productos
para aprovechar mejor el espacio (lo que en retail se conoce como \textit{space management}).

\section{¿A quién está dirigido este manual?}

A cualquier persona que vaya a \textbf{operar el panel}: cargar videos, calibrar las zonas de
análisis y leer los resultados. No se requieren conocimientos de programación. Si busca
documentación técnica sobre cómo está construido el sistema, consulte el \textbf{Manual
Técnico} que acompaña este documento.

\rvnota{Este sistema \textbf{no identifica personas}. Los rostros se difuminan automáticamente en
cada video analizado, antes de mostrarse en pantalla (ver capítulo \ref{cap:privacidad}).}

% ============================================================================
\chapter{Primeros pasos}
% ============================================================================

\section{Requisitos}

\begin{itemize}
  \item Un navegador web moderno (Chrome, Edge o Firefox recientes).
  \item Acceso a la dirección donde esté publicado el panel (por ejemplo,
    \texttt{http://localhost:5173} en una instalación local).
  \item Un video de la góndola a analizar, en formato \texttt{.mp4}, \texttt{.avi},
    \texttt{.mov}, \texttt{.mkv} o \texttt{.webm}, de hasta 500\,MB.
\end{itemize}

\section{La pantalla principal}

Al abrir el panel se ve una \textbf{barra lateral} a la izquierda con cinco secciones, y el
contenido de cada sección a la derecha:

\begin{center}
\begin{tabular}{>{\raggedright\arraybackslash}p{3.3cm} >{\raggedright\arraybackslash}p{10.2cm}}
\toprule
\rowcolor{rvCodeBg}\textbf{Sección} & \textbf{Para qué sirve} \\
\midrule
01 Lectura general & Resumen del último análisis: cifras clave y la escena analizada en vivo. \\
02 Analítica & Gráficas del comportamiento de los visitantes (embudo, permanencia, recorridos). \\
03 Plano de tienda & Mapa de calor y clasificación de cada góndola (caliente / templada / fría). \\
04 Evidencias & Lista de cada toma y devolución detectada, con su momento exacto. \\
05 Informe & Documento resumen, listo para imprimir o guardar como PDF. \\
\bottomrule
\end{tabular}
\end{center}

La barra lateral puede \textbf{contraerse} con el botón de flecha junto al logotipo, para dejar
más espacio a la pantalla.

% ============================================================================
\chapter{Flujo de trabajo: de un video a un resultado}
% ============================================================================

El uso típico del panel sigue cuatro pasos, siempre desde la sección \textbf{Lectura general}.

\section{\rvpaso{Paso 1 --- Subir el video}}

Haga clic en \textbf{``Seleccionar archivo de video''} y elija el video de la góndola desde su
computador. El sistema lo sube y, automáticamente, extrae un fotograma de \textbf{vista previa}
--- con los rostros ya difuminados --- para que pueda calibrar las zonas sin ver caras
reconocibles.

\rvnota{El análisis \textbf{todavía no se ejecuta} en este paso. Solo se prepara el video y la
vista previa.}

\section{\rvpaso{Paso 2 --- Calibrar las zonas}}

Con la vista previa en pantalla, haga clic en \textbf{``Editar zonas''}. Aparecerá un editor
donde cada góndola está marcada con dos polígonos:

\begin{itemize}
  \item Un polígono \textbf{grande}, que cubre el cuerpo de la persona (define cuándo alguien
    ``visita'' esa zona).
  \item Un polígono \textbf{punteado, más pequeño}, que marca dónde la \textbf{mano} toca el
    producto o el estante (define cuándo se cuenta una toma o devolución).
\end{itemize}

Arrastre los vértices (puntos) de cada polígono hasta que coincidan con los límites reales de la
góndola en la imagen. Un polígono que se cruza a sí mismo o queda sin área se marca en rojo y no
se puede guardar hasta corregirlo.

\rvconsejo{Si el video es muy parecido al anterior (misma cámara, mismo ángulo), normalmente no
hace falta recalibrar: las zonas guardadas se conservan entre análisis. Use \textbf{``Restablecer''}
dentro del editor solo si quiere volver a las zonas de fábrica.}

Al hacer clic en \textbf{``Guardar zonas''}, el sistema valida los polígonos y, si ya existe un
análisis previo del mismo video, lo \textbf{vuelve a ejecutar automáticamente} con las zonas
nuevas.

\section{\rvpaso{Paso 3 --- Ejecutar el análisis}}

Haga clic en \textbf{``Analizar video''}. El procesamiento corre en segundo plano; la barra de
estado muestra el progreso (\textit{``Preparando modelos y video''}, \textit{``Analizando el
video con YOLO''}\ldots). Según la duración del video y si el servidor cuenta con tarjeta
gráfica, esto puede tardar desde segundos hasta varios minutos.

\rvnota{No es necesario mantener la pestaña abierta todo el tiempo: al volver, el estado se
restaura automáticamente.}

\section{\rvpaso{Paso 4 --- Revisar los resultados}}

Cuando el estado cambia a \textbf{``completado''}, las cinco secciones del panel se llenan con
los datos del análisis. Puede pasar directamente a ``Analítica'', ``Plano de tienda'',
``Evidencias'' o ``Informe'' para explorarlos (capítulo \ref{cap:vistas}).

\section{Reiniciar métricas}

El botón \textbf{``Reiniciar métricas''} borra los resultados del análisis actual (no las zonas
calibradas, que se conservan) para empezar de cero con un nuevo video.

% ============================================================================
\chapter{Las cinco vistas del panel}
\label{cap:vistas}
% ============================================================================

\section{01 --- Lectura general}

Muestra, de un vistazo:

\begin{itemize}
  \item \textbf{Visitantes únicos}: personas distintas detectadas frente a la góndola.
  \item \textbf{Tomas detectadas}: productos retirados del estante.
  \item \textbf{Tasa de interacción}: qué porcentaje de visitantes llegó a tomar un producto.
  \item \textbf{Tasa de devolución}: de las tomas, cuántas terminaron en \textit{put-back}
    (se tomó el producto y se volvió a dejar).
  \item La \textbf{escena analizada}: una reproducción de los fotogramas ya anotados por el
    sistema (zonas dibujadas, etiquetas de cada persona), siempre con los rostros difuminados.
\end{itemize}

\section{02 --- Analítica}

Cuatro gráficas complementarias:

\begin{itemize}
  \item \textbf{Embudo de interacción}: cuántas personas pasan de ``tránsito'' a ``parada'',
    ``aproximación'', ``toma'' y, finalmente, cuántas no mostraron devolución.
  \item \textbf{Permanencia promedio}: cuánto tiempo se detiene la gente frente a cada góndola.
  \item \textbf{Comportamiento observado}: reparte a los visitantes entre ``sin toma'', ``toma
    sin devolución'' y ``toma y devolución''.
  \item \textbf{Flujo entre góndolas}: cómo se mueve la gente de una zona a otra (diagrama de
    flujo / Sankey).
\end{itemize}

\section{03 --- Plano de tienda}

Un \textbf{mapa de calor} muestra, con colores, dónde se concentró más la circulación y la
permanencia. A la derecha, cada góndola aparece con su \textbf{clasificación}:

\begin{center}
\begin{tabular}{lp{10.5cm}}
\toprule
\rowcolor{rvCodeBg}\textbf{Clasificación} & \textbf{Qué significa} \\
\midrule
\textcolor{rvRed}{Caliente} & Alto tráfico e interacción: el estante que mejor funciona hoy. \\
Templada & Tráfico intermedio: hay potencial de mejora. \\
\textcolor{rvMuted}{Fría} & Bajo tráfico: candidata a recibir productos de mayor rotación, o a
reducir el espacio que ocupa. \\
\bottomrule
\end{tabular}
\end{center}

Debajo del mapa, cada zona incluye una \textbf{recomendación en texto} generada automáticamente
a partir de sus propios datos (nunca es una opinión genérica: siempre cita las cifras observadas
de esa zona).

\section{04 --- Evidencias}

Una lista cronológica de cada \textbf{toma} y \textbf{devolución} detectada, con la hora dentro
del video, la zona donde ocurrió y un identificador temporal de la persona (que \textbf{no} es un
dato que permita reconocerla: se reinicia en cada análisis).

\section{05 --- Informe}

Un documento de una sola página, pensado para \textbf{imprimir o guardar como PDF}
(botón \textbf{``Imprimir / guardar PDF''}), con el resumen ejecutivo, los hallazgos principales,
el detalle por zona y las recomendaciones del agente. Útil para compartir con el equipo de
mercadeo o de operaciones sin necesidad de abrir el panel.

% ============================================================================
\chapter{Interpretando las recomendaciones}
% ============================================================================

Las recomendaciones del panel las genera un \textbf{motor de reglas} que solo usa los datos
observados por las cámaras (no es un modelo de lenguaje ni una opinión externa). Dos tipos de
recomendación aparecen en el panel:

\begin{itemize}
  \item \textbf{Por zona}: un diagnóstico y una acción sugerida (por ejemplo, ajustar precio o
    señalización si hay mucha ``consideración'' --- toma y devolución --- pero poca compra).
  \item \textbf{De planograma}: cruza la zona más caliente con la más fría y propone un reparto
    de espacio (\% de \textit{facing}) proporcional al interés observado en cada una.
\end{itemize}

\rvnota{Estas recomendaciones son un \textbf{apoyo a la decisión}, no un reemplazo del criterio
humano: deben contrastarse con datos de ventas, inventario y la realidad del punto de venta.}

% ============================================================================
\chapter{Accesibilidad}
% ============================================================================

El botón \textbf{``Accesibilidad''}, siempre visible en la esquina inferior derecha, abre un
panel de preferencias de visualización que se guardan en el propio dispositivo:

\begin{center}
\begin{tabular}{lp{10.5cm}}
\toprule
\rowcolor{rvCodeBg}\textbf{Opción} & \textbf{Efecto} \\
\midrule
Tamaño del texto & Ajusta toda la interfaz a 100\,\%, 112\,\% o 125\,\%. \\
Tipografía legible & Reduce los estilos decorativos de las fuentes. \\
Leer esta sección & Lee en voz alta el contenido de la pantalla actual (usa la voz del
dispositivo). \\
Alto contraste & Refuerza el contraste de textos, bordes y fondos. \\
Escala de grises & Elimina el color para reducir distracciones visuales. \\
Resaltar controles & Marca con un borde visible cada botón y enlace. \\
Cursor grande & Facilita apuntar y hacer clic. \\
Reducir movimiento & Desactiva animaciones y transiciones. \\
\bottomrule
\end{tabular}
\end{center}

Todo el panel puede recorrerse con el teclado (\texttt{Tab} para moverse, \texttt{Enter} para
activar, \texttt{Esc} para cerrar paneles).

% ============================================================================
\chapter{Privacidad: lo que el sistema nunca hace}
\label{cap:privacidad}
% ============================================================================

\begin{itemize}
  \item \textbf{No reconoce rostros.} No identifica ni guarda quién es cada persona.
  \item \textbf{Difumina las caras} de forma automática en todos los fotogramas que se muestran
    o se guardan, antes de cualquier otro procesamiento.
  \item \textbf{No envía el video a internet} para analizarlo: todo el procesamiento ocurre en
    el propio equipo donde está instalado el sistema.
  \item Los identificadores de cada persona en pantalla (``ID: 3'', por ejemplo) son
    \textbf{temporales}: se reinician en cada análisis y no se cruzan con ninguna otra base de
    datos.
\end{itemize}

Esto cumple con la \textbf{Ley 1581 de 2012} de Colombia sobre protección de datos personales.

% ============================================================================
\chapter{Preguntas frecuentes}
% ============================================================================

\begin{description}[leftmargin=0pt,style=nextline]
  \item[¿Por qué el video tarda en analizarse?] El tiempo depende de la duración del video y de
    si el equipo tiene tarjeta gráfica. Un video de un minuto suele tardar entre 15 segundos
    (con GPU) y unos pocos minutos (sin GPU).
  \item[¿Puedo subir otro video mientras uno se está analizando?] No: el sistema procesa un
    análisis a la vez y responderá que hay uno ``en curso'' hasta que termine.
  \item[¿Qué pasa si edito las zonas después de analizar?] El sistema relanza automáticamente el
    análisis sobre el mismo video con las zonas nuevas, sin que tenga que volver a subirlo.
  \item[¿Se pierde el análisis si cierro la pestaña?] No. El estado se guarda en el servidor y se
    restaura al volver a abrir el panel.
  \item[¿Las cifras en ``0'' significan un error?] No necesariamente: si aún no se ha subido ni
    analizado ningún video, el panel muestra ceros como estado inicial.
\end{description}

% ============================================================================
\chapter{Glosario}
% ============================================================================

\begin{description}[leftmargin=0pt,style=nextline]
  \item[Pick-up] Toma de un producto del estante por parte de un visitante.
  \item[Put-back] Devolución de un producto al estante tras haberlo tomado.
  \item[Dwell time] Tiempo de permanencia de una persona frente a una zona.
  \item[Heatmap (mapa de calor)] Representación visual de dónde se concentró más tráfico o
    permanencia.
  \item[Facing] Porcentaje de espacio de góndola asignado a una categoría de producto.
  \item[Planograma] Plan de distribución de productos sobre la góndola.
  \item[Zona de agarro] El área concreta, dentro de cada góndola, donde el sistema detecta el
    contacto de la mano con el producto.
\end{description}

% ============================================================================
\chapter{Soporte}
% ============================================================================

Para dudas sobre el uso del panel o para reportar un comportamiento inesperado, contacte al
equipo de RetailVision Edge AI a través del docente asesor del proyecto, \textbf{Sergio Arley
Puerto Moreno}, en la Universidad Santo Tomás, Seccional Tunja.

\end{document}
